\documentclass[11pt]{article}

\usepackage[margin=1.7cm]{geometry}
\usepackage{amsmath}
\usepackage{amssymb}
\usepackage{enumitem}
\usepackage{titlesec}
\usepackage{xcolor}
\usepackage{fancyhdr}
\usepackage{multicol}
\usepackage{array}

\titleformat{\section}{\large\bfseries\color{blue!50!black}}{\thesection}{0.6em}{}
\titlespacing{\section}{0pt}{1.1em}{0.4em}
\titleformat{\subsection}{\normalsize\bfseries}{\thesubsection}{0.6em}{}
\titlespacing{\subsection}{0pt}{0.8em}{0.3em}

\pagestyle{fancy}
\setlength{\headheight}{14pt}
\fancyhf{}
\fancyhead[L]{Factorising and Expanding Brackets}
\fancyhead[R]{Name: \makebox[4cm]{\hrulefill}}
\fancyfoot[C]{\thepage}
\renewcommand{\headrulewidth}{0.4pt}

\newcommand{\ans}{\hrulefill}
\newcommand{\anshere}{\par\vspace{0.2em}\noindent\ans\par\vspace{0.2em}}
\setlist[enumerate,1]{leftmargin=1.4em, itemsep=0.5em, topsep=0.2em}
\setlength{\columnsep}{1.1cm}
\setlength{\columnseprule}{0.4pt}

\begin{document}

\begin{center}
  {\Large\bfseries Factorising and Expanding Brackets}
\end{center}
\vspace{0.2em}
\noindent\textit{Four tiers: work through Fluency first, then Problem Solving and Reasoning
in either order, then try Enrichment as a challenge.}

\section*{Fluency}
\noindent\textit{Expand or factorise each expression. One bracket, no application.}

\begin{multicols}{2}
\subsection*{Expand}
\begin{enumerate}[label=(\alph*)]
  \item $3(x+4)$
  \anshere
  \item $5(2x-3)$
  \anshere
  \item $-2(x+6)$
  \anshere
  \item $4(3x-5)$
  \anshere
  \item $-(x-7)$
  \anshere
  \item $x(x+3)$
  \anshere
  \item $2x(x-5)$
  \anshere
  \item $-3x(2x+1)$
  \anshere
\end{enumerate}

\columnbreak

\subsection*{Factorise}
\begin{enumerate}[label=(\alph*)]
  \item $4x+12$
  \anshere
  \item $6x-15$
  \anshere
  \item $x^2+5x$
  \anshere
  \item $3x^2-9x$
  \anshere
  \item $10x+25$
  \anshere
  \item $2x^2+8x$
  \anshere
  \item $9x-6$
  \anshere
  \item $9x^2-6x$
  \anshere
\end{enumerate}
\end{multicols}

\newpage

\section*{Problem Solving}
\noindent\textit{These need more than one step: multiple brackets, variable common
factors, or an equation to solve.}

\begin{multicols}{2}
\subsection*{Expand and simplify}
\begin{enumerate}[label=(\alph*)]
  \item $3(x+2)+2(x-5)$
  \anshere
  \item $4(2x-1)-3(x+2)$
  \anshere
  \item $(x+3)(x+4)$
  \anshere
  \item $(2x-1)(x+5)$
  \anshere
  \item $(x-6)(x-2)$
  \anshere
\end{enumerate}

\subsection*{Factorise fully}
\begin{enumerate}[label=(\alph*)]
  \item $x^2+7x+10$
  \anshere
  \item $x^2-3x-18$
  \anshere
  \item $x^2-16$
  \anshere
  \item $4x^2-9$
  \anshere
  \item $2x^2+x-6$
  \anshere
\end{enumerate}

\columnbreak

\subsection*{Solve}
\noindent\textit{Expand the brackets first, then solve for $x$.}
\begin{enumerate}[label=(\alph*)]
  \item $3(x+5)=27$
  \anshere
  \item $2(3x-4)=5x+1$
  \anshere
  \item $5(x-2)-3(x+1)=7$
  \anshere
\end{enumerate}

\vspace{1em}
\subsection*{A step further}
\begin{enumerate}[label=(\alph*)]
  \item A rectangle has length $(x+5)$ and width $(x-2)$. Write an expression for its
    area, expanded and simplified.
  \anshere
  \item The length of a rectangle is $2x-3$ and its width is $x$. The area is
    $2x^2-3x$. Factorise this area expression to check it matches the length and
    width given.
  \anshere
\end{enumerate}
\end{multicols}

\newpage

\section*{Reasoning}
\noindent\textit{No variables here --- use expanding or factorising as a mental-maths
trick to work these out \emph{without} a calculator. Show the expanded/factorised form
you used.}

\begin{enumerate}[label=(\alph*)]
  \item Write $99$ as $(100-1)$ and $101$ as $(100+1)$. Use the difference of two
    squares to calculate $99 \times 101$.
  \anshere
  \item Use factorising to quickly calculate $37 \times 8 + 37 \times 2$.
  \anshere
  \item Write $52$ as $(50+2)$ and $48$ as $(50-2)$. Use the difference of two squares
    to calculate $52 \times 48$.
  \anshere
  \item Write $203$ as $(200+3)$. Use expanding to calculate $6 \times 203$.
  \anshere
  \item A school has $29$ classes with $31$ students in each. Write $29$ as $(30-1)$
    and $31$ as $(30+1)$, then use the difference of two squares to find the total
    number of students.
  \anshere
  \item Write $999$ as $(1000-1)$. Use expanding to calculate $999^2$.
  \anshere
\end{enumerate}

\section*{Enrichment}
\noindent\textit{Inequalities behave like equations, except multiplying or dividing by
a negative number reverses the inequality sign.}

\begin{enumerate}[label=(\alph*)]
  \item Solve: \quad $3(x+2) > 15$
  \anshere
  \item Solve: \quad $2(x-5) \le 4(x+1)$
  \anshere
  \item Solve: \quad $-3(x+1) < 2(x-4)$
  \anshere
  \item Solve the double inequality: \quad $2 < 3(x-1) \le 11$
  \anshere
  \item A taxi charges a \$4 call-out fee plus \$3 per kilometre travelled. Write an
    inequality for the total cost of a trip of $k$ kilometres, then solve it to find
    how far you could travel and keep the cost under \$40.
  \anshere
\end{enumerate}

\end{document}
